\documentclass[11pt]{article}
\usepackage[margin=1in]{geometry}
\usepackage{enumitem}
% =============================================================================
% Preambles/header.tex — shared LaTeX/Beamer preamble for this project
%
% Use from a Beamer lecture by prepending:
%     \input{header}
% and compiling with TEXINPUTS=../Preambles:$TEXINPUTS (the /compile-latex
% skill does this automatically).
%
% PALETTE CONTRACT. Color names defined here MUST match the names in
% Quarto/theme-template.scss so Beamer and Quarto renderings use the same
% palette. The scripts/check-palette-sync.sh script enforces this.
%
% To customize: edit the HEX values below AND the matching $variable in
% Quarto/theme-template.scss, then run scripts/check-palette-sync.sh.
% =============================================================================

% -----------------------------------------------------------------------------
% Engine detection + encoding
%
% Our /compile-latex skill uses XeLaTeX, where inputenc is unnecessary and
% can produce encoding warnings. Only load inputenc under pdfTeX.
% -----------------------------------------------------------------------------
\usepackage{iftex}
\ifPDFTeX
  \usepackage[utf8]{inputenc}
\fi

\usepackage{amsmath, amssymb, amsthm}
\usepackage{booktabs}
\usepackage{graphicx}

% -----------------------------------------------------------------------------
% PALETTE — mirrors Quarto/theme-template.scss variable names.
% Load xcolor and define colors BEFORE anything that references them
% (notably hyperref, hypersetup, and the Beamer color assignments).
% -----------------------------------------------------------------------------
\usepackage{xcolor}

\definecolor{primary-blue}{HTML}{012169}      % headings, accents
\definecolor{primary-gold}{HTML}{B9975B}      % emphasis, borders
\definecolor{highlight-yellow}{HTML}{F2A900}  % markers, alerts
\definecolor{light-bg}{HTML}{E8EDF5}          % title / section backgrounds
\definecolor{jet}{HTML}{1A1A1A}               % body text

% Semantic colors (used in diagrams and annotations)
\definecolor{positive}{HTML}{15803D}          % good / observed / identified
\definecolor{negative}{HTML}{B91C1C}          % bad / problematic / bias
\definecolor{neutral}{HTML}{525252}           % reference / context

% Highlight accents (match .hi-* classes in the SCSS)
\definecolor{hi-slate}{HTML}{314F4F}
\definecolor{hi-green}{HTML}{2E8B57}
\definecolor{hi-red}{HTML}{C41E3A}

% Semantic aliases so skill templates and snippets can write
% \textcolor{accent}{...} without hard-coding "primary-blue".
\colorlet{accent}{primary-blue}
\colorlet{accent2}{primary-gold}

% -----------------------------------------------------------------------------
% Hyperref — loaded AFTER xcolor + palette so link colors resolve.
% -----------------------------------------------------------------------------
\usepackage{hyperref}
\hypersetup{colorlinks=true, linkcolor=primary-blue, urlcolor=primary-blue,
            citecolor=primary-blue, pdfborder={0 0 0}}

% -----------------------------------------------------------------------------
% Course code — set once per deck (after \input{header}, before \begin{document})
% via \coursecode{CS401}. Shown in the footer on every slide so a deck's course
% is identifiable without opening the title page. Empty by default.
% -----------------------------------------------------------------------------
\makeatletter
\newcommand{\coursecode}[1]{\gdef\@coursecodetext{#1}}
\gdef\@coursecodetext{}
\makeatother

% -----------------------------------------------------------------------------
% Beamer theme assignments (only apply under Beamer).
%
% We detect Beamer via \@ifclassloaded{beamer}, which is the documented,
% non-internal way. This must be wrapped in \makeatletter/\makeatother
% because \@ifclassloaded contains an @-catcode character.
% -----------------------------------------------------------------------------
\makeatletter
\@ifclassloaded{beamer}{%
  \usefonttheme{professionalfonts}
  \setbeamercolor{structure}{fg=primary-blue}
  \setbeamercolor{frametitle}{fg=primary-blue}
  \setbeamercolor{title}{fg=primary-blue}
  \setbeamercolor{subtitle}{fg=primary-gold}
  \setbeamercolor{author}{fg=primary-blue}
  \setbeamercolor{institute}{fg=neutral}
  \setbeamercolor{date}{fg=primary-gold}
  \setbeamercolor{itemize item}{fg=primary-blue}
  \setbeamercolor{itemize subitem}{fg=primary-gold}
  \setbeamercolor{alerted text}{fg=primary-gold}
  \setbeamercolor{block title}{fg=white, bg=primary-blue}
  \setbeamercolor{block body}{bg=light-bg}
  % Semantic block variants: block=definition/concept (blue), exampleblock=
  % worked example (green/positive), alertblock=key takeaway or caution (gold).
  % Keep to <=2 colored boxes per slide across ALL three variants (INV-7).
  \setbeamercolor{block title example}{fg=white, bg=positive}
  \setbeamercolor{block body example}{bg=light-bg}
  \setbeamercolor{block title alerted}{fg=white, bg=primary-gold}
  \setbeamercolor{block body alerted}{bg=light-bg}

  % Minimal footer: course code (left) + page X / N (right), no navigation clutter
  \setbeamertemplate{navigation symbols}{}
  \setbeamertemplate{footline}{%
    \color{neutral}\footnotesize\hspace{0.7em}\@coursecodetext\hfill
    \insertframenumber/\inserttotalframenumber\hspace{0.7em}\vspace{0.3em}}
}{}
\makeatother

% -----------------------------------------------------------------------------
% TikZ setup — libraries the snippet gallery uses
% -----------------------------------------------------------------------------
\usepackage{tikz}
\usetikzlibrary{arrows.meta, positioning, calc, decorations.pathreplacing,
                fit, shapes.geometric, backgrounds}

% Shared TikZ styles — snippets and hand-written diagrams can pick these up
% via `\begin{tikzpicture}[every node/.style={...}]` or by naming specific
% styles (e.g., `\node[dag-node] (X) at (0,0) {X};`).
\tikzset{
  % Node styles for causal DAGs and similar
  dag-node/.style = {
    draw=primary-blue, fill=light-bg, rounded corners=2pt,
    minimum width=1.2cm, minimum height=0.9cm, align=center, font=\small
  },
  decision-node/.style = {
    draw=primary-gold, fill=white, diamond, aspect=1.8,
    minimum width=1.8cm, minimum height=1.2cm, align=center, font=\small,
    inner sep=1pt
  },
  % Edge styles — observed vs counterfactual
  observed-edge/.style = {-{Latex[length=2.5mm]}, thick, draw=jet},
  counterfactual-edge/.style = {-{Latex[length=2.5mm]}, thick, draw=neutral, dashed},
  confound-edge/.style = {-{Latex[length=2.5mm]}, thick, draw=negative, dashed},
  % Data-point styles for plots
  observed-dot/.style = {fill=primary-blue, draw=primary-blue, circle, inner sep=1.2pt},
  counterfactual-dot/.style = {fill=white, draw=neutral, circle, inner sep=1.2pt}
}

% -----------------------------------------------------------------------------
% Convenience macros
% -----------------------------------------------------------------------------
\newcommand{\muted}[1]{\textcolor{neutral}{#1}}
\newcommand{\key}[1]{\textcolor{primary-gold}{\textbf{#1}}}
\newcommand{\good}[1]{\textcolor{positive}{#1}}
\newcommand{\bad}[1]{\textcolor{negative}{#1}}

% Standout frame macro: full-bleed dark background for section transitions.
% Beamer-only (uses \begin{frame}/\setbeamercolor/\usebeamerfont) -- do not
% call from an article-class file (e.g. Notes/<CODE>/*-notes.tex). \muted,
% \key, \good, \bad, and \coursecode above are plain xcolor/gdef and are
% safe to use from either class; verified when Notes/article-class support
% was added.
% Expand: \transitionslide{Your transition title}
\newcommand{\transitionslide}[1]{%
  \begingroup
  \setbeamercolor{background canvas}{bg=primary-blue}
  \begin{frame}[plain]
    \centering
    \vfill
    {\usebeamerfont{title}\color{white}\Large #1\par}
    \vfill
  \end{frame}
  \endgroup
}

% Section divider frame (Beamer-only), an alternative to \transitionslide with a
% darker two-line style: near-black (jet) background, a small white label line
% above a larger gold title, and a thin gold rule along the bottom edge.
% Expand: \sectiondivider{Section 2}{Pointers}
\newcommand{\sectiondivider}[2]{%
  \begingroup
  \setbeamercolor{background canvas}{bg=jet}
  \begin{frame}[plain]
    \vspace*{3.0cm}
    {\color{white}\ttfamily\large #1\par}
    \vspace{0.5cm}
    {\color{primary-gold}\LARGE #2\par}
    \begin{tikzpicture}[remember picture, overlay]
      \draw[primary-gold, line width=2.5pt]
        ($(current page.south west)+(0,0.1)$) -- ($(current page.south east)+(0,0.1)$);
    \end{tikzpicture}
  \end{frame}
  \endgroup
}

% -----------------------------------------------------------------------------
% Channel watermark — "Concepts That Click" (YouTube).
%
% Article-class files (CompetitiveExam Q/A sets, Notes, Assignments) get a
% light diagonal draft watermark on every page plus a centered footer naming
% the channel. Beamer decks get a subtle TikZ background watermark instead
% (the footline already carries the course code). Centralized here so every
% MCQ PDF is branded without per-file edits.
% -----------------------------------------------------------------------------
\makeatletter
\@ifclassloaded{article}{%
  \usepackage{draftwatermark}
  \SetWatermarkText{Concepts That Click}
  \SetWatermarkScale{1}
  \SetWatermarkFontSize{2cm}
  \SetWatermarkLightness{0.86}
  \SetWatermarkAngle{45}
  \usepackage{fancyhdr}
  \pagestyle{fancy}
  \fancyhf{}
  \fancyfoot[C]{\footnotesize\textcolor{neutral}{Concepts That Click~|~YouTube: \href{https://www.youtube.com/@concepts.thatclick}{@concepts.thatclick}}}
  \renewcommand{\headrulewidth}{0pt}
  \renewcommand{\footrulewidth}{0pt}
  \fancypagestyle{plain}{%
    \fancyhf{}%
    \fancyfoot[C]{\footnotesize\textcolor{neutral}{Concepts That Click~|~YouTube: \href{https://www.youtube.com/@concepts.thatclick}{@concepts.thatclick}}}%
    \renewcommand{\headrulewidth}{0pt}%
    \renewcommand{\footrulewidth}{0pt}%
  }
}{}
\@ifclassloaded{beamer}{%
  \setbeamertemplate{background}{%
    \begin{tikzpicture}[remember picture, overlay]
      \node[rotate=45, opacity=0.055, align=center]
        at (current page.center)
        {\Huge\textcolor{neutral}{Concepts That Click}};
    \end{tikzpicture}%
  }
}{}
\makeatother 
\coursecode{JKSSB}
\renewcommand{\thesection}{64.\arabic{section}}
\newtheorem{example}{Example}[section]
\newenvironment{definitionbox}[1]{%
  \par\noindent\textbf{Definition (#1).}\ \itshape
}{\par\vspace{0.3em}}
\title{Current Cybersecurity, Cryptography and Cyber Law --- Detailed Lecture Notes}
\author{JKSSB Computer Awareness}
\date{\today}

\begin{document}
\maketitle

\section{Scope and study method}
This lesson covers the current-affairs edge of the official Computer Awareness
scope: computer virus and latest anti-virus, terms and abbreviations used in
IT, and the role of IT in governance. It treats cybersecurity threats,
cryptography basics, firewall types, and cyber law at awareness level. The
exam depth is recall and classification, with at most a one-step scenario.
The study method for every term is the near-neighbour contrast: what it is,
what it is not, and the nearest confusion listed alongside it.

\begin{center}
\begin{tabular}{p{0.24\textwidth}p{0.66\textwidth}}
\toprule
\textbf{Term} & \textbf{Meaning in this lesson} \\
\midrule
Malware & Hostile software, including virus, worm, Trojan horse, ransomware, spyware. \\
Parasitic / file-infector virus & A virus that attaches itself to an executable file. \\
Phishing & Fraudulent messages that trick the user into revealing credentials. \\
Ransomware & Malicious code that locks or encrypts files and demands payment. \\
DDoS & Distributed Denial of Service: flood of traffic from many machines. \\
Hacking & Unauthorised access to a computer or account. \\
DES --- Data Encryption Standard & Symmetric cipher working on 64-bit blocks. \\
Encryption & Conversion of plaintext to ciphertext for confidentiality. \\
Digital signature & Authenticity and integrity proof; not confidentiality. \\
Firewall & Barrier between a trusted and an untrusted network. \\
Personal data & Data identifying a person, protected by consent and purpose limits. \\
Copyright & Owner's exclusive rights over an original work. \\
\bottomrule
\end{tabular}
\end{center}

\section{Malware and the parasitic virus}
\textbf{Malware} is the general name for hostile software. A \textbf{virus}
is malware that needs a host program and some human action, such as opening
an attachment, to spread. A \textbf{worm} spreads on its own across a
network without a host program. A \textbf{Trojan horse} looks like useful
software but hides a harmful payload, and it does not self-replicate.

\begin{definitionbox}{Parasitic (file-infector) virus}
A virus that attaches itself to an executable file. When the infected program
runs, the virus code runs first and then hands control back to the host
program, spreading further as infected files are shared or copied.
\end{definitionbox}

The parasitic virus is the file-infecting member of the family: it lives
inside executable files, activates when the host program runs, and travels
with shared files. An \textbf{anti-virus} detects, blocks, and removes such
infections, and its database must be kept updated against the latest strains.

\begin{example}
A student copies a free game from a friend's pen drive and runs it on the
office desktop. The game file carries a file infector: the virus runs first,
attaches to other executable files, and spreads when those files are shared
further. Updating the anti-virus and scanning the drive before opening the
files is the correct defence.
\end{example}

\section{Phishing, ransomware, DDoS, and hacking}
These four threats differ in mechanism, and the exam swaps their definitions.
\textbf{Phishing} attacks the user: a fraudulent mail or message, often styled
as a bank or KYC notice, tricks the user into revealing a password or OTP on
a fake login page. \textbf{Ransomware} attacks the data: malicious code
locks or encrypts the victim's files and demands payment for the decryption
key. \textbf{DDoS} (Distributed Denial of Service) attacks availability: a
flood of traffic from many machines overwhelms a server so genuine users
cannot reach it. A flood from a single machine is only DoS; the distributed
prefix matters. \textbf{Hacking} is the broad term for unauthorised access to
a computer or account; phishing, ransomware, and DDoS are specific methods or
goals within that broad idea.

\begin{definitionbox}{Ransomware}
Malicious code that locks or encrypts the victim's files and demands payment
for their release.
\end{definitionbox}

\begin{example}
An item describes a clerk whose files are suddenly unreadable, with a message
demanding payment for the key. Tracing the mechanism --- files locked,
payment demanded --- gives ransomware. If the stem instead described a fake
bank mail asking for an OTP, the same trace gives phishing; a server slowed to
a halt under flood traffic gives DDoS.
\end{example}

\section{Cryptography and DES}
\textbf{Cryptography} protects information by transforming it so only the
intended recipient can read it. \textbf{Plaintext} is the readable message,
\textbf{ciphertext} is the transformed message, and a \textbf{key} controls
the transformation. \textbf{Encryption} converts plaintext to ciphertext and
\textbf{decryption} converts it back. In \textbf{symmetric} cryptography the
same key encrypts and decrypts; in \textbf{asymmetric} cryptography a public
key and a private key form a pair.

\begin{definitionbox}{DES (Data Encryption Standard)}
A symmetric cipher that processes data in 64-bit blocks. Its key is 64 bits
long, of which 56 bits are effective, the remaining 8 bits being parity.
\end{definitionbox}

The two numbers for DES must be kept apart. The \textbf{block size} is 64
bits: that is the chunk of data processed at a time. The \textbf{effective
key} is 56 bits. An option offering 128 bits confuses DES with a later
cipher's block or key size.

\section{Encryption vs digital signature}
This is the TRAP-15 distinction and it recurs in papers. \textbf{Encryption}
provides \textbf{confidentiality}: only the intended recipient can read the
message. A \textbf{digital signature} provides \textbf{authenticity and
integrity}: it proves who sent the message and that it was not altered in
transit. A signature does not hide the message content, and encryption by
itself does not prove the sender's identity.

\begin{definitionbox}{Digital signature}
A cryptographic proof attached to a message that establishes the sender's
authenticity and the message's integrity, without by itself providing
confidentiality.
\end{definitionbox}

\begin{example}
An item states, ``A digital signature encrypts the message body so only the
recipient can read it.'' This is wrong. Hiding the body is encryption
(confidentiality). The signature answers two different questions: who sent it
and whether it was changed. The wrong option has merged the two security
verbs.
\end{example}

\section{Firewalls: three types}
A \textbf{firewall} is a barrier between a trusted internal network and an
untrusted external network. It permits or denies traffic according to a rule
set. It filters packets and connections; it does not disinfect
file-infecting viruses inside downloads that its rules allow through --- that
is the anti-virus's job.

A \textbf{packet-filtering} firewall examines each packet's headers ---
source and destination address, port, and protocol --- and decides packet by
packet. It is the fastest and simplest type but cannot judge a session as a
whole. A \textbf{stateful-inspection} firewall examines the same headers plus
the \textbf{connection state}: it remembers open sessions and drops packets
that do not belong to a known session. An \textbf{application (proxy)}
firewall sits between the client and the server as a go-between and inspects
application-layer content and commands, such as web or mail requests.

\begin{definitionbox}{Firewall}
A network barrier that permits or denies traffic between a trusted internal
network and an untrusted external network according to a rule set.
\end{definitionbox}

The ordering shorthand is: header only, then header plus state, then full
content via proxy. Filtering checks the envelope, stateful inspection checks
the conversation, and the proxy reads the letter inside.

\section{Internet governance, privacy, and data protection}
The Internet has no single owner; names and numbers are coordinated
internationally, for example by ICANN (Internet Corporation for Assigned
Names and Numbers). Governance concerns at awareness level are four: the
\textbf{security} of networks, the \textbf{privacy} of personal data, fair
\textbf{access} to the network, and which country's \textbf{law} applies when
an act crosses borders. The positive side is the \textbf{role of IT in
governance} (e-governance): online services, transparency, and delivery of
benefits to citizens.

\textbf{Personal data} is any data that identifies a person, such as a name
combined with a phone number, biometrics, or account details. Protection
principles at awareness level are: collect data with \textbf{consent}, use it
only for the stated \textbf{purpose}, keep it \textbf{secure}, and retain it
no longer than needed. Privacy settings, strong passwords, and OTP discipline
are the user-side defences. Misuse or leakage of personal data can attract
legal consequences under cyber law.

\section{Cyber law and copyright: awareness-level map}
Cyber law at this level means two things: unauthorised access, fraud, and
data misuse attract punishment, and creators hold enforceable rights over
their works. The \textbf{Copyright Act} protects original literary, artistic,
musical, and software works from copying without permission. The
awareness-level map matches each concept to its area: the \textbf{meaning of
copyright} (what exclusive rights the owner holds), \textbf{first ownership}
(who owns the work first), \textbf{infringement} (what counts as violation),
\textbf{fair dealing} (limited permitted uses such as private study), and
\textbf{offences} (punishment for wilful infringement). Software piracy ---
copying licensed software without permission --- is the standard infringement
example. Section numbers are recall only where the question itself supplies
them; otherwise the concept match is what the exam tests.

\section{Exam retrieval and common confusions}
Seven distinctions prevent most errors on this topic.
\begin{enumerate}
  \item DES block size is \textbf{64 bits}; the effective key is 56 bits. Do
    not offer 128 for either (DES fact).
  \item Digital signature $\ne$ encryption: authenticity/integrity vs
    confidentiality (TRAP-15/TRM-11).
  \item A file infector attaches to \textbf{executable files}; it is not a
    worm (self-spreading) or a Trojan (non-replicating disguise).
  \item Packet filter checks headers only; stateful adds \textbf{connection
    state}; proxy inspects \textbf{application content}.
  \item Phishing steals \textbf{credentials}; ransomware locks \textbf{files};
    DDoS blocks \textbf{access}; hacking is unauthorised access generally.
  \item DDoS needs \textbf{many} machines (distributed); a single-machine
    flood is only DoS.
  \item Copyright protects the owner against \textbf{infringement}, with
    limited \textbf{fair-dealing} exceptions such as study --- piracy is the
    standard violation example.
\end{enumerate}

\begin{example}
An item asks which firewall remembers open sessions and drops stray packets.
Trace the triple: the packet filter decides per packet, the proxy inspects
application content, and only stateful inspection tracks connection state.
The answer is the stateful-inspection firewall. If the stem instead asked
which device disinfects an infected executable, the answer leaves firewalls
entirely and gives the anti-virus.
\end{example}

\subsection*{Exercises}
\begin{enumerate}[label=\arabic*.]
  \item Define a parasitic (file-infector) virus and distinguish it from a worm and a Trojan horse.
  \item State the DES block size and effective key size, and expand the abbreviation DES.
  \item Distinguish packet-filtering, stateful-inspection, and application/proxy firewalls.
  \item Distinguish phishing, ransomware, DDoS, and hacking in one line each.
  \item Explain why a digital signature is not encryption, using TRAP-15.
  \item List the four Internet-governance concerns and the four data-protection principles from this lesson.
  \item Match each copyright concept (meaning, first ownership, infringement, fair dealing, offences) to what it covers.
\end{enumerate}

\subsection*{Solutions}
\begingroup\small
\begin{enumerate}[label=\arabic*.]
  \item A parasitic virus attaches to an executable file, runs when the host
    program runs, and spreads with shared files. A worm spreads on its own
    across a network without a host; a Trojan horse disguises itself as useful
    software and does not self-replicate.
  \item DES (Data Encryption Standard) is a symmetric cipher with a 64-bit
    block size and a 56-bit effective key (64 bits with parity).
  \item Packet-filtering checks per-packet headers only; stateful inspection
    adds connection-state tracking; the application/proxy firewall inspects
    application-layer content as a go-between.
  \item Phishing tricks the user into revealing credentials; ransomware locks
    or encrypts files for ransom; DDoS floods a service from many machines so
    genuine users cannot reach it; hacking is unauthorised access generally.
  \item Encryption provides confidentiality (who can read it); a digital
    signature provides authenticity and integrity (who sent it, whether it was
    changed). Merging them is TRAP-15.
  \item Governance concerns: security, privacy, access, jurisdiction
    (applicable law). Data-protection principles: consent, stated purpose,
    security, limited retention.
  \item Meaning: the exclusive rights the owner holds; first ownership: who
    owns the work first; infringement: what counts as violation (for example
    software piracy); fair dealing: limited permitted uses such as private
    study; offences: punishment for wilful infringement.
\end{enumerate}
\endgroup
\end{document}
